\documentclass[12pt,letterpaper,oneside]{article}
\input{ITESCA/maestria-en-gestion-administrativa/plan-de-negocios-mga/formato-itesca-plan-negocios-generadas.tex}
\begin{document}
\portadaITESCA{Diagrama de flujo de producción o servicio}{6550}
\section{Diagrama de flujo de producción o servicio}

Este documento constituye un \textbf{borrador hipotético pendiente de confirmación} para AM Taller Autocentro. Presenta una propuesta del proceso de recepción, diagnóstico y atención de solicitudes relacionadas con mecánica, llantas y refacciones. La secuencia deberá validarse mediante observación directa de la operación y confirmación de tiempos, recursos, responsables y controles. No representa evidencia de servicios realizados, inventario disponible, acuerdos con proveedores ni operación comercial validada.

El propósito es relacionar las actividades del servicio con sus decisiones, tiempos, recursos y responsables. Las dos versiones siguientes describen el mismo proceso: la primera muestra los nodos y sus conexiones; la segunda detalla cada paso mediante una tabla operativa.

\subsection{Versión gráfica del proceso propuesto}

Los nodos de actividad representan operaciones y los identificados como decisión muestran rutas condicionales. La numeración vincula esta representación con la versión redactada.

\begin{tabularx}{\textwidth}{|c|X|X|}
\hline
\textbf{Nodo} & \textbf{Actividad o decisión} & \textbf{Conexión} \\
\hline
1 & \textbf{Inicio: recepción de la solicitud} & Continúa al nodo 2. \\
\hline
2 & \textbf{Registro de datos y requerimiento} & Continúa al nodo 3. \\
\hline
3 & \textbf{Decisión: información suficiente} & Sí: nodo 4. No: regresar al nodo 2 para solicitar aclaraciones. \\
\hline
4 & \textbf{Programación de cita o revisión} & Continúa al nodo 5. \\
\hline
5 & \textbf{Recepción e inspección inicial} & Continúa al nodo 6. \\
\hline
6 & \textbf{Diagnóstico o verificación técnica} & Continúa al nodo 7. \\
\hline
7 & \textbf{Decisión: se requieren piezas, llantas o insumos} & Sí: nodo 8. No: nodo 9. \\
\hline
8 & \textbf{Consulta de inventario propuesto y proveedores} & Recurso localizable: nodo 9. No localizable: nodo 11. \\
\hline
9 & \textbf{Elaboración y comunicación de la propuesta} & Continúa al nodo 10. \\
\hline
10 & \textbf{Decisión: existe autorización} & Sí: nodo 12. No: nodo 11. \\
\hline
11 & \textbf{Cierre sin ejecución o reprogramación} & Continúa al nodo 16. \\
\hline
12 & \textbf{Preparación del área, herramientas e insumos} & Continúa al nodo 13. \\
\hline
13 & \textbf{Ejecución del servicio autorizado} & Continúa al nodo 14. \\
\hline
14 & \textbf{Revisión de calidad y seguridad} & Conforme: nodo 15. No conforme: regresar al nodo 13. \\
\hline
15 & \textbf{Explicación del trabajo y devolución propuesta} & Continúa al nodo 16. \\
\hline
16 & \textbf{Registro de cierre y seguimiento} & Fin del proceso. \\
\hline
\end{tabularx}

El flujo incorpora controles antes de cualquier intervención. Se requiere confirmar el requerimiento, comunicar la propuesta de servicio y obtener autorización. Las referencias consultadas en catálogos se consideran \textbf{por cotizar}; por ello, no deben interpretarse como productos existentes, compatibles, disponibles o listos para venta.

\subsection{Versión redactada del proceso}

Los tiempos indicados son \textbf{supuestos preliminares de planeación}, no mediciones obtenidas en campo. El diagnóstico, la reparación y la obtención de insumos pueden requerir periodos diferentes según el vehículo, el alcance autorizado y la complejidad técnica.

\begin{longtable}{|p{0.22\textwidth}|p{0.12\textwidth}|p{0.26\textwidth}|p{0.16\textwidth}|}
\hline
\textbf{Actividad detallada} & \textbf{Tiempo propuesto} & \textbf{Materiales, equipo o herramientas} & \textbf{Personal requerido} \\
\hline
\endfirsthead
\hline
\textbf{Actividad detallada} & \textbf{Tiempo propuesto} & \textbf{Materiales, equipo o herramientas} & \textbf{Personal requerido} \\
\hline
\endhead
Recibir la solicitud e identificar si corresponde a mecánica, llantas, refacciones o revisión general. & 5 a 10 minutos & Formato de recepción, agenda y medio de comunicación. & Responsable de recepción. \\
\hline
Registrar datos de contacto, características generales del vehículo, necesidad expresada y disponibilidad. No afirmar compatibilidades sin verificación. & 5 a 10 minutos & Formato de registro y referencias por cotizar. & Responsable de recepción. \\
\hline
Comprobar que la información permita programar la atención. Si faltan datos, solicitar las aclaraciones necesarias. & 3 a 5 minutos & Registro y lista de datos mínimos. & Responsable de recepción. \\
\hline
Proponer fecha y horario sujetos a la capacidad operativa que posteriormente se confirme. & 5 minutos & Agenda de citas y datos de contacto. & Recepción o coordinación. \\
\hline
Recibir el vehículo y documentar sus condiciones observables y el servicio solicitado. & 10 a 20 minutos & Formato de recepción, equipo de seguridad e instrumentos básicos. & Personal técnico y recepción. \\
\hline
Realizar el diagnóstico dentro del alcance autorizado y registrar los hallazgos sin ampliar automáticamente el servicio. & 20 a 90 minutos & Herramientas de diagnóstico y equipo de protección. & Personal técnico competente. \\
\hline
Determinar si se requieren refacciones, llantas u otros insumos y verificar sus especificaciones antes de proponer referencias. & 10 a 20 minutos & Fichas técnicas, catálogos y registro de requerimientos. & Personal técnico y responsable de inventario propuesto. \\
\hline
Consultar la disponibilidad interna, si llegara a existir, o solicitar cotizaciones a proveedores potenciales. & Variable & Control de inventario propuesto, directorio por validar y cotizaciones. & Compras o coordinación. \\
\hline
Integrar una propuesta con actividades, insumos, tiempo estimado y condiciones, y comunicarla antes de intervenir. & 10 a 20 minutos & Diagnóstico, cotizaciones y formato de autorización. & Personal técnico y recepción. \\
\hline
Obtener una autorización verificable. Si no se concede, cerrar el registro o plantear una posible reprogramación. & 5 a 15 minutos & Medio de autorización aceptado. & Persona solicitante y recepción. \\
\hline
Preparar el área y comprobar que las herramientas e insumos correspondan al trabajo autorizado. & 10 a 30 minutos & Equipo de protección, herramientas e insumos verificados. & Personal técnico. \\
\hline
Ejecutar únicamente las actividades autorizadas. Una necesidad adicional exige detener el proceso y solicitar nueva confirmación. & Variable & Herramientas y equipo específicos según el servicio. & Personal técnico competente. \\
\hline
Revisar la calidad y seguridad. Toda desviación debe documentarse y corregirse antes de plantear la devolución del vehículo. & 15 a 30 minutos & Lista de verificación e instrumentos aplicables. & Personal técnico y responsable de revisión. \\
\hline
Explicar las actividades efectuadas, recomendaciones y limitaciones sin prometer resultados no comprobados. & 10 a 15 minutos & Registro del servicio y recomendaciones documentadas. & Personal técnico y recepción. \\
\hline
Cerrar el registro, resguardar la documentación y acordar seguimiento cuando resulte pertinente. & 5 a 10 minutos & Expediente del servicio y agenda. & Responsable de recepción. \\
\hline
\end{longtable}

\subsection{Criterios de control}

\begin{itemize}
\item Confirmar que cada solicitud cuente con información suficiente antes de agendarla.
\item No asumir existencia o compatibilidad de llantas, refacciones o insumos.
\item Presentar el diagnóstico y la propuesta antes de solicitar autorización.
\item Evitar la ejecución de actividades adicionales sin confirmación.
\item Registrar la revisión de calidad y seguridad antes de la devolución.
\item Comparar posteriormente los tiempos reales con los rangos propuestos.
\end{itemize}

Para convertir este borrador en un trabajo real se requieren los tiempos medidos de cada actividad, los puestos y el número de responsables, las herramientas disponibles, la capacidad de atención, los criterios técnicos de revisión, el procedimiento de autorización, las reglas de resguardo de datos, el funcionamiento real de la agenda, el método de control de inventario y los proveedores efectivamente evaluados. También permanecen pendientes de confirmación las extensiones admitidas para los archivos de ambas versiones, la rúbrica de evaluación y el peso individual asignado a cada entrega.
\clearpage
\printbibliography[title={Fuentes de consulta}]
\end{document}
